\documentclass{beamer}

\usetheme{metropolis}
\usecolortheme{dolphin}

\title{Fast Answers, Occasional Lies}
\author{Siddharth \& Amit}
\date{\today}

\newcommand{\Z}{\mathbb{Z}}
\newcommand{\bld}[1]{\textbf{#1}}
\newcommand{\itl}[1]{\textit{#1}}
\newcommand{\N}{\mathbb{N}}
\newcommand{\Q}{\mathbb{Q}}
\newcommand{\R}{\mathbb{R}}
\newcommand{\eps}{\varepsilon}
\newcommand{\poly}{\mathrm{poly}}
\newcommand{\polylog}{\mathrm{polylog}}
\newcommand{\bigO}{\mathcal{O}}

\begin{document}

\maketitle

%------------------------------------------------
\section{Heuristic Primality Testing}

\begin{frame}{Why Heuristics?}

\begin{itemize}
    \item Exact methods (e.g., trial division) are inefficient
    \item We seek \bld{fast} methods for large integers
    \item Trade-off: speed vs certainty
\end{itemize}

\end{frame}

%------------------------------------------------

\begin{frame}{General Strategy}

\begin{itemize}
    \item Use algebraic properties satisfied by primes
    \item Check if $n$ behaves like a prime
    \item Reject quickly if a contradiction appears
\end{itemize}


\begin{block}{Key Idea}
Primes satisfy strong structural constraints
\end{block}

\end{frame}

%------------------------------------------------
\section{Euclid's Algorithm}

\begin{frame}{The GCD Primitive}

\[
\gcd(a,b) = \gcd(b, a \bmod b)
\]


\begin{itemize}
    \item Based on the division algorithm
    \item Reduces problem size at each step
\end{itemize}

\end{frame}

%------------------------------------------------

\begin{frame}{Algorithm}

Given $a > b$:

\begin{enumerate}
    \item $r = a \bmod b$
    \item $a \leftarrow b,\quad b \leftarrow r$
    \item Repeat until $b = 0$
    \item Output $a$
\end{enumerate}

\end{frame}

%------------------------------------------------

\begin{frame}{Example}

\[
\gcd(48,18)
= \gcd(18,12)
= \gcd(12,6)
= \gcd(6,0)
= 6
\]

\end{frame}

%------------------------------------------------

\begin{frame}{Why It Matters}

\begin{itemize}
    \item Runs in $\bigO(\log n)$ time

    \item Detects non-trivial factors efficiently
\end{itemize}


\begin{block}{Application}
If $\gcd(a,n) > 1$, then $n$ is composite
\end{block}

\end{frame}

%------------------------------------------------
\section{Wilson's Theorem}

\begin{frame}{Statement}

\begin{block}{Wilson's Theorem}
$n$ is prime $\iff (n-1)! \equiv -1 \bmod n$
\end{block}

\end{frame}

\begin{frame}{Proof}
Suppose that $n$ is composite. Then there exists a prime $q$ such that
\[
2 \leq q < n \quad \text{and} \quad q \mid n.
\]
Hence, there exists an integer $k$ such that
\[
n = qk.
\]
\pause
Assume for contradiction that
\[
(n-1)! \equiv -1 \pmod{n}.
\]
Then there exists an integer $m$ such that
\[
(n-1)! = nm - 1 = (qk)m - 1 = q(km) - 1.
\]

\end{frame}

\begin{frame}{Proof}
Thus, $(n-1)!$ is one less than a multiple of $q$, which implies
\[
(n-1)! \equiv -1 \pmod{q}.
\]
\pause
On the other hand, since $2 \leq q \leq n-1$, the integer $q$ appears among the factors in
\[
(n-1)! = (n-1)(n-2)\cdots 2 \cdot 1.
\]
Therefore,
\[
(n-1)! \equiv 0 \pmod{q}.
\]

This is a contradiction. Hence,
$
(n-1)! \not\equiv -1 \pmod{n}
$
when $n$ is composite.
    

\end{frame}

%------------------------------------------------

\begin{frame}{Interpretation}

\begin{itemize}
    \item Complete characterization of primes

    \item No false positives or negatives
\end{itemize}


\begin{block}{Philosophy}
Primes exhibit strong symmetry in modular arithmetic
\end{block}

\end{frame}

%------------------------------------------------

\begin{frame}{Why Not Use It?}

\begin{itemize}
    \item Requires computing $(n-1)!$

    \item Time complexity: $\Omega(n)$
\end{itemize}


\begin{block}{Conclusion}
Exact, but computationally infeasible
\end{block}

\end{frame}

\section{Selfridge's Conjecture}

\begin{frame}{Statement}

Let $p$ be odd with $p \equiv \pm 2 \bmod 5$


If:
\[
2^{p-1} \equiv 1 \bmod p
\]
\[
F_{p+1} \equiv 0 \bmod p
\]


Then $p$ is prime

\end{frame}

%------------------------------------------------

\begin{frame}{What is Happening?}

\begin{itemize}
    \item First condition: Fermat test (base 2)

    \item Second condition: Fibonacci structure
\end{itemize}


\begin{block}{Insight}
Combining independent constraints strengthens the test
\end{block}

\end{frame}

%------------------------------------------------

\begin{frame}{Generalization}

\begin{itemize}
    \item Uses Fibonacci polynomials $F_n(x)$

    \item Based on quadratic non-residues
\end{itemize}


\begin{block}{Theme}
Arithmetic + algebraic structure $\Rightarrow$ stronger tests
\end{block}

\end{frame}

%------------------------------------------------

\begin{frame}{Status}

\begin{itemize}
    \item Selfridge, Pomerance and Wagstaff together offered \$620 for a counterexample or proof that there is none, with the prize now due from the Number Theory Foundation.
    \item Still unproven
\end{itemize}


\begin{block}{Impact}
Led to the Baillie--PSW primality test
\end{block}

\end{frame}

%------------------------------------------------

\begin{frame}{Summary}

\begin{itemize}
    \item Euclid: fast arithmetic backbone

    \item Wilson: exact but impractical

    \item Fermat: fast but imperfect

    \item Selfridge: stronger heuristic, still unproven
\end{itemize}

\end{frame}

\end{document}